\scanpage{18}
\S\,1 : \uline{Les foncteurs} $H^i_{\unsure{\mathscr C|\tau|\mathfrak C}}$

\noindent\makebox[0pt][r]{\smash{\raisebox{-12\baselineskip}{\(\left\{\rule{0pt}{24\baselineskip}\right.\)}}\hspace{0.6em}}%
1. Sur les foncteurs $A$-linéaires (de modules et modules topologiques)

2. Définition des $H^i_{\mathscr C}(M)$, premières propriétés.

3. Fonctorialité pour les $H^i_{\mathscr C}(M)$ \inserted{(y inclure \unsure{cup-produit|cup-produits})}

4. Cas d'un anneau local $A$ : \inserted{premières propriétés spéciales; \uline{caractérisation des} $H^{\unsure{i}}$}

\deleted{\unsure{4}}5. \inserted{relations avec la codim. coh.} Rappels sur les Modules de Cohen-Macaulay

6. Résultats auxiliaires sur certains $\operatorname{Ext}_A^i(N,M)$ \margin{[$M$ \unsure{Ch. Mac|Coh. Mac.} de supp. $\operatorname{Spec}(A)$}

\deleted{\unsure{6}}7. Cas particuliers

\S\,2 \uline{Foncteurs et modules dualisants} \inserted{(théorie locale)} \uline{Le th. de dualité locale}.

1. Définition et premières propriétés des foncteurs et modules dualisants.

2. Extension du foncteur dualisant

3. Th. d'existence et unicité; exemples.

4. \deleted{\unsure{Le|Les|La} \unsure{théorème|théorèmes|théorie} de dualité locale} Structure des $H^i(M)$.

5. Module fondamental et \uline{th. de dualité locale}

6. Anneaux de Gorenstein.


\S\,3 \uline{Les Foncteurs} $H_i$, \uline{et les complexes dualisants}
% archive sheet number: 17

